\documentclass[12pt]{article}

\usepackage[margin=1.5cm]{geometry}
\usepackage[colorlinks]{hyperref}
\usepackage{natbib}

\title{Vision Statement}
\author{Joaquín Rapela}

\begin{document}

\maketitle

\section{Research Vision}

\textit{Mathematics without biology or biology without mathematics are not as
wonderful}.

Since pursuing my PhD, my research career has been highly interdisciplinary
spanning analytical (e.g., mathematics, signal processing, statistics, machine
learning, computer science) and neuroscience experimental (e.g., statistical
analysis of responses of single cells, electrocorticography, Utah arrays,
Neuropixels recordings, close-loop experimental control) subjects. I delight
when I create a new method to process biological data, but I delight much
further when I can apply this method to discover new aspects of brain function.

My \emph{Statistical Neuroscience Research Group} will initially focus on
creating advanced statistical methods to characterise neural and behavioral
data, offline and in real time, for short- and long-duration experiments. In
addition, in collaboration with experimental groups, we will apply these
methods, or those created by others, to investigate behavioural and neural
problems.

\subsection{Latent variable models}

It is challenging to extract meaning from the large number of spikes generated
by high-channel count recordings probes, e.g., Neuropixels. To address this
problem, my supervisor at the Gatsby Unit, Prof.~Maneesh Sahani, and
collaborators, created Gaussian Process Factor
Analysis~\citep[GPFA][]{yuEtAl09}, a probabilistic method to extract
low-dimensional representation (i.e., latent variables) from high-dimensional
electrophysiological recordings. At the Unit I have distributed a Python
implementation of a recent extension of GPFA, Sparse Variational Gaussian Process
Factor Analysis~\citep[svGPFA][]{dunckerAndSahani18}, which uses a sparse
variational method to substantially reduces the memory and time complexity
for the estimation of GPFA models and, critically, allows point-process
observations, eliminating the need of binning spikes. I have improved and
optimized svGPFA, and explored its capabilities to
\href{https://joacorapela.github.io/svGPFA/auto_examples/index.html}{characterize
several neural datasets}.

In collaboration with the laboratory of Marcus Stephenson-Jones at the SWC, I
have extracted svGPFA neural latents from striatal Neuropixels recordings in
mice performing a stereotyped poking sequence task~\citep{thompsonEtAl24}. I
then used these neural latents as inputs to a Hidden Markov Model to decode,
with impressive accuracy, the mice behavioral states. Critically, this is the
first position decoder that uses the spiking activity of whole populations
and does not require the estimation of place fields of single units.

We are now observing that patterns of neural latents observed during behavior
replay during sleep periods. If these observations are confirmed, it will be
the first time that replay is identified using continuous latent variables of
neural population activity, bypassing the estimation of place fields, which
could open a new chapter in memory research.

In collaboration with Prof.~Maneessh Sahani, my group will continue developing
svGPFA and will collaborate with experimental groups at NPP and CDB to
investigate the performance of svGPFA for position decoding and replay
detection in the hippocampus.

\subsection{Naturalistic, long-duration and continual experimentation}

At the SWC and Gatsby Unit, we are pioneering a new type of naturalistic,
long-duration and continual experimentation, the Aeon project, that will
revolutionise neuroscience experimentation~\citep{campagnerEtAl25}. 

In 2023-2024 I mentored a very talented undergraduate student from NPP,
\href{https://zimoli02.github.io/}{Ms.~Zimo Li}, in a
\href{https://www.simonsfoundation.org/2023/09/26/simons-foundation-announces-next-class-of-neuroscience-undergraduate-research-fellows/}{Simons
Foundation SURFiN
fellowship}, where she performed the first comprehensive statistical analysis
of Aeon data.

The achievements of Zimo were outstanding. She processed very complex data,
that at that stage were poorly organised, and skillfully applied sophisticated
statistical models (e.g., linear dynamical systems, hidden Markov models, point
process regression) and produced the first statistical Aeon discoveries, which
were included in the
\href{https://www.biorxiv.org/content/10.1101/2025.07.31.664513v1}{Aeon
platform paper}, with Zimo as coauthor. The collaboration with Zimo taught me
that I can do excellent research with students that do not have a solid
mathematical background.

In collaboration with the Aeon team at the Gatsby Unit and SWC, my group will
continue developing statistical methods for the analysis of Aeon data. Due to
the large size of Aeon recordings, conventional batch machine learning (ML) methods
cannot be used to process Aeon data. Aeon requires online methods able to
process non-stationary data, and are computationally very efficient to handle
very large datasets. A detailed description of our ML data
processing plans for Aeon appear in a
\href{https://www.gatsby.ucl.ac.uk/~rapela/npp/submittedProposal.pdf}{collaborative
grant proposal between the Gatsby Unit and the Allen Institute for Neural
Dynamics we submitted to BBSRC/NSF in 2024}.

\subsection{Intelligent experimental control}

Current neuroscience experimentation is mostly performed in open loop with
predefined sets of rules designed by experimenters. More informative
experiments can be designed if they could be controlled in close loop by
intelligent inferences performed on behaviour and neural recordings.

Bonsai is today the most widely used software for close loop experimental
control. In 2022 I realised that adding ML functionality to
Bonsai could make possible a radically new type of close loop intelligent
experimentation, and lead the successful application of a BBSRC proposal
(\href{https://gtr.ukri.org/projects?ref=BB\%2FW019132\%2F1}{BB/W019132/1}) to
build the \href{https://bonsai-rx.org/machinelearning/}{Bonsai.ML package}.

Most current ML methods operate offline, trained and performing inference on
datasets previously recorded. However, Bonsai.ML requires a different type of ML
methods, that are trained and perform inference on streams continuously
delivering data in real time. We build this type of ML into Bonsai.ML. It now
contains probabilistic ML methods for real time processing of behavioral (e.g.,
kinematics inference with linear dynamical systems, behavioral state estimation
with hidden Markov models) and neural data (neural latent variables inference
with linear dynamical systems, hyppocampal position decoding and replay detection with a
point process state-space model). Please refer to the
\href{https://bonsai-rx.org/machinelearning/}{documentation of
Bonsai.ML}.

We are not integration large language models (LLMs) into Bonsai. First, we are
using LLMs to help Bonsai users build correct and optimised workflows using
natural language. Second, we are adding an LLM node to Bonsai to incroporate
the large semantic knowledge of LLMs into the experimental controls loop.
Please refer to the \href{https://github.com/joacorapela/os4ls}{Open Source for
the Life Science proposal} we are currently preparing to fund these
developments.

My group will continue developing online probabilistic models to enable
intelligent experimental control in Bonsai.

\subsection{Collaborations}

\section{Teaching Vision}

\bibliographystyle{apalike}
\bibliography{latentVariableModels,replay}

\end{document}
